% Auto-typeset from the previously authored formula-by-formula study manuscript.
% Based on the thesis main branch f80e5ac903dd, 2026-10-31.

\chapter*{数学预备知识}
\addcontentsline{toc}{chapter}{数学预备知识}

关联原文：第二至第五章。建议第一次学习时先掌握前七节，再随着正文返回查阅。\par

\section{变量、复数与矩阵的记号}

本研究的空间谱问题以 x 为微分变量，以 z 为演化变量；物理上 x 是推迟时间，但在 IST 中习惯称 x 方程为空间 Lax 方程。谱参数 \ensuremath{\lambda} 是额外引入的\textbf{复变量}，不是 x 或 z 的函数。\par

如果 f 是复数，\textbackslash{}overline f 是共轭，|f|\textsuperscript{2}=f\textbackslash{}overline f；矩阵 B 的共轭转置为 B\textasciicircum{}\textbackslash{}dagger=\textbackslash{}overline B\textasciicircum{}\{\textbackslash{}mathsf T\}。常用 Pauli 矩阵：\par

\[
\sigma_1=\begin{pmatrix}0&1\\1&0\end{pmatrix},
\qquad
\sigma_3=\begin{pmatrix}1&0\\0&-1\end{pmatrix}.
\]

直接相乘得\par

\[
\sigma_3\sigma_1=
\begin{pmatrix}0&1\\-1&0\end{pmatrix},
\qquad
\sigma_1\sigma_3=-\sigma_3\sigma_1.
\]

“反对易”即两个乘法次序结果相反。若 B=\textbackslash{}begin\{pmatrix\}0\&u\textbackslash{}\textbackslash{}\textbackslash{}overline u\&0\textbackslash{}end\{pmatrix\}，也有 B\textbackslash{}sigma\_3=-\textbackslash{}sigma\_3 B。定义交换子 [X,Y]=XY-YX；非对角矩阵 X=\textbackslash{}begin\{pmatrix\}0\&a\textbackslash{}\textbackslash{}b\&0\textbackslash{}end\{pmatrix\} 满足\par

\[
[\sigma_3,X]=\begin{pmatrix}0&2a\\-2b&0\end{pmatrix}.
\]

所以与 \ensuremath{\sigma}\textsubscript{3} 做交换子，恰能抽取势矩阵的非对角分量；这就是第四章无穷远重构公式的来源。\par

\section{Lax 对、相容性和零曲率是怎么来的}

假设一个两分量辅助向量或基本解矩阵满足\par

\[
\Phi_x=U(x,z,\lambda)\Phi,
\qquad
\Phi_z=V(x,z,\lambda)\Phi.
\]

这是同一个函数沿 x、z 两个方向的线性方程。假定可交换求偏导：\par

\[
\begin{aligned}
\Phi_{xz}&=U_z\Phi+U\Phi_z=(U_z+UV)\Phi,\\
\Phi_{zx}&=V_x\Phi+V\Phi_x=(V_x+VU)\Phi.
\end{aligned}
\]

相减：\par

\[
(U_z-V_x+UV-VU)\Phi=0.
\]

若 \ensuremath{\Phi} 是可逆基本解，就得到\par

\[
U_z-V_x+[U,V]=0.
\]

这个矩阵恒等式对所有允许 \ensuremath{\lambda} 成立。若 U、V 是 \ensuremath{\lambda} 的多项式或有理式，则可把同次幂、同一极点处的留数逐一比较。所得多个矩阵分量方程，正是 HMB 的非线性方程。\par

\textbf{理解要点}：Lax 对不是另一个无关物理模型，而是把非线性 PDE 的三个未知场编码到两个线性谱方程里。我们要通过谱问题“换一条路”求出非线性方程解。\par

\section{什么是背景谱与 Jost 解}

从方程的背景 u\ensuremath{\to}Q 得到常系数矩阵 U\ensuremath{\infty}(\ensuremath{\lambda})。常系数 ODE\par

\[
\psi_x=U_\infty(\lambda)\psi
\]

可以通过求特征值和特征向量精确求解。非零背景的常系数解一般是两种振荡模式的线性组合。\par

真实场 u(x,z) 只要求 x\ensuremath{\to}\ensuremath{\pm}\ensuremath{\infty} 时趋于 Q，中间允许有局域扰动。因此从左、右两个无穷远分别构造“看起来像背景振荡模式”的解 \ensuremath{\chi}\textsubscript{-}、\ensuremath{\chi}\textsubscript{+}，称为 Jost 解。它们解的是\textbf{同一个}空间 ODE，只是归一化端点不同。\par

把它类比常微分方程的两套基本解：同一个解空间的两套基之间存在一个可逆常矩阵 S。这个矩阵在谱问题中叫散射矩阵。S 的非对角元衡量模式之间的混合，透射与反射系数便从它的分量定义。\par

不要误认为 Jost 解与 RHP 矩阵 M 相同：M 是按照解析性，把左右 Jost 的合适列重新拼接、必要时除以透射系数后得到的对象。\par

\section{变分常数如何形成 Volterra 方程}

先从标量\par

\[
y'(x)=ay(x)+f(x)
\]

开始。令 y=e\textasciicircum{}\{ax\}c(x)，得到 c'(x)=e\textasciicircum{}\{-ax\}f(x)。积分：\par

\[
y(x)=e^{a(x-x_0)}y(x_0)+
\int_{x_0}^{x}e^{a(x-s)}f(s)\,\mathrm{d}s.
\]

对于矩阵 A，指数变成 e\textasciicircum{}\{A(x-s)\}，公式结构不变。若我们知道 y 在 x\ensuremath{\to}-\ensuremath{\infty} 与背景解一致，就令 x\textsubscript{0}\ensuremath{\to}-\ensuremath{\infty}；若在 +\ensuremath{\infty} 一致，积分方向从 +\ensuremath{\infty} 到 x。积分方程中未知函数只在积分号内以一次形式出现，属于 Volterra 型。\par

它为何有用？在适合的谱区域中，指数核有界，连续迭代得到\par

\[
y=y_0+\mathcal K y_0+\mathcal K^2 y_0+\cdots.
\]

第 n 次迭代是有序 n 重积分，产生 1/n! 因子。这通常使 Volterra 方程比任意 Fredholm 方程更容易证明唯一可解。\par

\section{为什么矩阵解只有部分列解析}

形式上含 e\textasciicircum{}\{\textbackslash{}pm 2\textbackslash{}mathrm\{i\}\textbackslash{}rho(x-y)\} 的积分核，\ensuremath{\rho} 为复数。写 \ensuremath{\rho}=a+\textbackslash{}mathrm\{i\}b，则\par

\[
\left|e^{2\mathrm{i}\rho(x-y)}\right|
=e^{-2b(x-y)}.
\]

当 x-y\ensuremath{\ge}0 且 b>0，指数不会增长；如果积分方向相反，就需要另一种符号。不同的端点（\ensuremath{\pm}\ensuremath{\infty}）和不同的列，产生四种搭配。\par

所以论文第三章的结论并不是“整个 \ensuremath{\chi}\textsubscript{-} 在 D\textsubscript{+} 解析”；实际是“\ensuremath{\chi}\textsubscript{-} 的某一列及 \ensuremath{\chi}\textsubscript{+} 的另一列在 D\textsubscript{+} 解析”。理解这种选列规则，后面 RHP 的两侧矩阵定义就不再是死记。\par

\section{双叶谱曲面为何出现}

对非零常背景，谱根是\par

\[
\rho^2=\lambda^2+Q^2.
\]

在 \ensuremath{\lambda}=\ensuremath{\pm}\textbackslash{}mathrm\{i\}Q 处平方根为零。固定 \ensuremath{\lambda} 时通常有两个可能值 \ensuremath{\rho}、\ensuremath{-}\ensuremath{\rho}，因而无法在整个 \ensuremath{\lambda} 平面选一个无歧义的解析分支；这导致双叶黎曼曲面。\par

令\par

\[
w=\lambda+\rho,
\qquad
\lambda=\dfrac{1}{2}\left(w-\dfrac{Q^2}{w}\right),
\qquad
\rho=\dfrac{1}{2}\left(w+\dfrac{Q^2}{w}\right).
\]

现在 \ensuremath{\lambda}、\ensuremath{\rho} 都成为 w 的单值有理函数，整个双叶问题被改写为 w 平面的问题。换叶对应\par

\[
h(w)=-\dfrac{Q^2}{w},
\qquad
\rho(h(w))=-\rho(w).
\]

背景连续谱变成实轴和圆 |w|=Q。\textbf{反射与离散谱的“成对、四点”对称，都与这个换叶映射密切相关。}\par

\section{什么是 RHP？为什么有跳跃和留数？}

最简单的矩阵 Riemann–Hilbert 问题是：在复平面一条定向轮廓 \ensuremath{\Sigma} 的两侧寻找分片解析矩阵 M\textsubscript{+}、M\textsubscript{-}，满足\par

\[
M_+(w)=M_-(w)J(w),\qquad w\in\Sigma,
\]

并满足指定的无穷远行为。J 叫跳跃矩阵。\par

散射关系能够告诉你“某条曲线两侧的两套解析列之间差什么”，所以从 Jost 解自然得到 J。\par

如果透射系数 s\textsubscript{1}\textsubscript{1}(\ensuremath{\zeta})=0 且为简单零点，除法 1/s\textsubscript{1}\textsubscript{1} 会产生极点。复分析给出\par

\[
\operatorname*{Res}_{w=\zeta}\dfrac{f(w)}{s_{11}(w)}
=\dfrac{f(\zeta)}{s'_{11}(\zeta)}.
\]

因此 RHP 不仅有跳跃，还必须指定离散点的留数。单值化谱对称会使一个自由谱点通常带出四个相关点。\par

\textbf{无反射}是指反射系数在整个连续谱上为零，而不是 q=0，也不是没有离散谱。此时 J=I；若轮廓上没有不可去奇性，M 可以表示为背景项加有限个简单极点项。求解 RHP 就变成有限维线性代数。\par

\section{留数条件与部分分式如何产生线性系统}

以一个列为例，设\par

\[
\operatorname*{Res}_{w=\zeta}M_1(w)=C\,M_2(\zeta),
\qquad M_2 \text{ 在 }\zeta\text{ 解析}.
\]

那么 M 的奇异主部就是\par

\[
\dfrac{C M_2(\zeta)e_1^{\mathsf T}}{w-\zeta}.
\]

这里 e\textsubscript{1}\textsuperscript{T} 是行向量 (1,0)，右乘的效果是“只把这一极点项放入第一列”。如果有多个离散点，就把主部逐个求和。\par

未知量不是任意函数，而是有限个正则列值 M\textsubscript{2}(\ensuremath{\zeta}\textsubscript{j}) 等。将部分分式在这些点取值，就得到有限个线性方程；这便是第四、五章的重要计算思路。\par

\section{无穷远展开、Liouville 定理和重构}

假设\par

\[
M(w)=I+\dfrac{M^{(1)}}{w}+O(w^{-2}),\qquad w\to\infty.
\]

把它放入空间 Lax 关系，对比 w\textsuperscript{0} 或 \ensuremath{\lambda}\textsuperscript{0} 的系数，会得到势矩阵 A(u) 与 [\ensuremath{\sigma}\textsubscript{3},M\textsuperscript{(}\textsuperscript{1}\textsuperscript{)}] 的关系。这叫“无穷远重构”。\par

Liouville 定理的用法：一个在整个复平面解析、且无穷远趋于零的矩阵函数，其每个分量都是有界整函数，因此恒为零。无反射部分分式推导和第五章恢复 Lax 对都在使用这种“排除了所有极点后只剩下多项式/常数”的论证。\par

\section{为什么 Maxwell–Bloch 变量需要另一种重构？}

HMB 比普通 NLS 多了 p、\ensuremath{\eta}。它们被放入传播矩阵的有理项\par

\[
\dfrac{\mathrm{i}R}{\lambda-a},
\qquad R=\begin{pmatrix}\eta&-p\\-\overline p&-\eta\end{pmatrix}.
\]

光场对应空间 Lax 矩阵在谱无穷远的非对角系数；介质矩阵却对应传播 Lax 矩阵在 \ensuremath{\lambda}=a 的\textbf{留数}。\par

介质方程还可以等价写成\par

\[
R_x=[U(a),R].
\]

如果 \ensuremath{\chi}(x,a) 是可逆空间基本解，则\par

\[
\dfrac{\mathrm{d}}{\mathrm{d}x}
\bigl(\chi^{-1}R\chi\bigr)=0.
\]

因此 R=\ensuremath{\chi} C \ensuremath{\chi}\textsuperscript{-}\textsuperscript{1}，C 为不依赖 x 的矩阵。由适当的背景与共振点正则性可确定 C。这个机制\textbf{不需要预先用 Darboux 构造孤子}，但必须处理好共振极点处的相容条件。\par

\section{易误解的逻辑关系}

\begin{itemize}[leftmargin=2em]
\item 零曲率给出 PDE：前提是 Lax 对与所需正则性已经建立。
\item 直接散射给出一个 RHP：不等于对任意抽象反射数据都证明了逆问题存在。
\item 有限线性系统可逆：不等于谱对称、分支点条件、背景端态自然全部满足。
\item 数值残差很小：验证计算实现，不代替严格证明。
\item 一个四点轨道的 RHP 解存在：不自动推广为任意 N 个轨道、任意复谱参数都可解。
\end{itemize}

\section{学习自检}

\begin{enumerate}[leftmargin=2em]
\item 用矩阵乘法亲自证明 U\_z\ensuremath{-}V\_x+[U,V]=0 的来源。
\item 用 \ensuremath{\rho}\textsuperscript{2}=\ensuremath{\lambda}\textsuperscript{2}+Q\textsuperscript{2} 推出 \ensuremath{\lambda}(w) 与 \ensuremath{\rho}(w)。
\item 对 x\ensuremath{-}y 的两个符号分别判断 |e\textasciicircum{}\{2\textbackslash{}mathrm\{i\}\ensuremath{\rho}(x\ensuremath{-}y)\}| 是否有界。
\item 用一次 Taylor 展开证明简单零点的留数公式。
\item 解释为什么 r\ensuremath{\equiv}0 仍可有非平凡孤子。
\item 证明 \ensuremath{\chi}\textsuperscript{-}\textsuperscript{1}R\ensuremath{\chi} 在 x 方向常数的公式，不跳过求导。
\item 用一句话区别本论文的“直接 RHP 求解”与“先 Darboux 后核对 RHP”。
\end{enumerate}

如果其中有一步仍不熟悉，可以先只学习相应小节，不必急着读完整论文证明。\par

